\subsection{Лабораторная работа № 2}

\subsubsection{Задача}

Необходимо реализовать таблицу
\mintinline{text}{Typeface_format_weight_stats}, содержимое которой повторяет
результат представления, созданного в первой лабораторной работе. Таблица
должна хранить для каждого начертания идентификатор гарнитуры, название
начертания, число поддерживаемых форматов и количество градаций насыщенности.

Актуальность агрегированных данных необходимо обеспечить с помощью триггеров
на таблицах \mintinline{text}{Font_family} и
\mintinline{text}{Typeface}. При изменении исходных данных соответствующие
строки таблицы статистики должны обновляться автоматически.

Требуется реализовать пять триггеров, покрывающих следующие изменения:

\begin{enumerate}
    \item изменение идентификатора гарнитуры в таблице
    \mintinline{text}{Font_family};
    \item удаление гарнитуры из таблицы \mintinline{text}{Font_family};
    \item добавление начертания в таблицу \mintinline{text}{Typeface};
    \item изменение гарнитуры, названия, формата или насыщенности начертания в
    таблице \mintinline{text}{Typeface};
    \item удаление начертания из таблицы \mintinline{text}{Typeface}.
\end{enumerate}

В рамках работы была подготовлена обновлённая схема базы данных: в неё
добавлена таблица статистики, возле таблиц \mintinline{text}{Font_family} и
\mintinline{text}{Typeface} указаны соответствующие триггеры, а их воздействие
на агрегированные данные показано пунктирными стрелками.

\subsubsection{Схема базы данных}

В исходную схему добавлена таблица
\mintinline{text}{Typeface_format_weight_stats}. Пунктирные стрелки показывают
операции, при выполнении которых триггеры изменяют содержащиеся в ней
агрегированные значения. Обновлённая схема представлена на
рисунке~\ref{fig:db-schema-triggers}.

\fullPageFigure{tex/graphics/db-schema-triggers.pdf}
{Схема базы данных с таблицей статистики и триггерами}
{\label{fig:db-schema-triggers}}
